\section{Tangent spaces and slice rank}\label{sec:tan}

\subsection{Algebraic geometry essentials}

We will need only a very small number of basic concepts from algebraic geometry, which we quickly review next. All the material here can be found in standard textbooks (e.g.,~\cite{Harris,Shafarevich}).
A \emph{variety} $\V{V}$ is the set of solutions, in an algebraically closed field, of some finite set of polynomials.
More formally, for a field $\FF$,\footnote{We henceforth denote by $\overline{\FF}$ the algebraic closure of the field $\FF$.}
the variety $\V{V} \sub \overline{\FF}^n$ \emph{cut out}
by the polynomials $f_1,\ldots,f_m \in \FF[x_1,\ldots,x_n]$ is 
$$\V{V} = \VV(f_1,\ldots,f_m) := \{ \B{x} \in \overline{\FF}^n \mid  f_1(\B{x})=\cdots=f_m(\B{x})=0\} .$$
We say that $\V{V} \sub \overline{\FF}^n$ is \emph{defined} over $\FF$ as it can be cut out by polynomials whose coefficients lie in $\FF$.
The ideal of $\V{V}$ is $\I(\V{V})=\{f \in \FF[\B{x}] \mid \forall p \in \V{V} \colon f(p)=0\}$.
Any variety $\V{V}$ can be uniquely written as the union of \emph{irreducible} varieties, where a variety is said to be irreducible if it cannot be written as the union of strictly contained varieties.
The \emph{dimension} of a variety $\V{V}$, denoted $\dim\V{V}$, is the maximal length $d$ of a chain of irreducible varieties 
$\emptyset \neq \V{V}_1 \subsetneq\cdots\subsetneq \V{V}_d \subsetneq \V{V}$.
The \emph{codimension} of $\V{V} \sub \overline{\FF}^n$ is simply $\codim\V{V}=n-\dim\V{V}$.


\subsection{Notation}\label{subsec:notation}

In the rest of the paper we will often find it convenient to identify, with a slight abuse of notation,
a bilinear map $f \colon \FF^{n_1} \times \FF^{n_2} \to \FF^m$ (or tensor) with a linear subspace of matrices, or \emph{matrix space}, $\V{L} \preceq \FF^{n_1 \times n_2}$.
If $f=(f_1,\ldots,f_m)$, we will identify $f$ with the linear subspace $\V{L}$ spanned by the $m$ matrices corresponding to the bilinear forms $f_1,\ldots,f_m$.
Note that this identification is not a correspondence, as it involves choosing a basis for $\V{L}$.  
Importantly, however, since the notions of tensor rank that we study are invariant under the action of the general linear group $\mathrm{GL}_n$ on each of the axes, the choice of basis we make is immaterial in the definition of rank, meaning that $\GR(\V{L})$, $\SR(\V{L})$, $\AR(\V{L})$ are nevertheless well defined. 

For the reader's convenience, we summarize below the different perspectives of tensor/bilinear map/matrix space that we use, and how they relate to each other:
\begin{itemize}
	\item A tensor $T = (a_{i,j,k}) \in \FF^{n_1\times n_2 \times n_3}$,  
	or a multilinear form $T(\B{x},\B{y},\B{z}) = \sum_{i,j,k} a_{i,j,k}x_iy_jz_k$.
	\item A bilinear form $f=(f_1,\ldots,f_{n_3}) \colon \FF^{n_1} \times \FF^{n_2} \to \FF^{n_3}$ with $f_k(\B{x},\B{y}) = \sum_{i,j} a_{i,j,k}x_iy_j$.
	\item A matrix space $\V{L} \preceq \FF^{n_1 \times n_2}$ spanned by $\{A_1,\ldots,A_{n_3}\}$ where $A_k = (a_{i,j,k})_{i,j}$.
\end{itemize}



\subsection{Tangent spaces of a variety}

For a variety $\V{V} \sub \KK^n$, the tangent space $\T_p\V{V}$ to $\V{V}$ at the point $p \in \V{V}$ is the linear subspace
$$\T_p\V{V} = \Big\{ \B{v} \in \KK^n \,\Big\vert\, \forall g \in \I(\V{V}) \colon \frac{\partial g}{\partial \B{v}}(p) = 0 \Big\}.$$
Equivalently, for any choice of a generating set $\{g_1,\ldots,g_s\} \sub \KK[x_1,\ldots,x_n]$ for the ideal $\I(\V{V})$ (which is finitely generated by Hilbert's basis theorem), the tangent space at $p \in \V{V}$ is $\T_p\V{V} = \ker \J_p$, where $\J_p$ is the Jacobian matrix
$$\begin{pmatrix}
	\frac{\partial{g_1}}{\partial x_1}(p) & \cdots & \frac{\partial{g_1}}{\partial x_n}(p) \\
	\vdots  & \ddots & \vdots  \\
	\frac{\partial{g_s}}{\partial x_1}(p) & \cdots & \frac{\partial{g_s}}{\partial x_1}(p) 
\end{pmatrix}_{s\times n} .$$

We will need the following basic fact about tangent spaces (for a proof see, e.g., Theorem~2.3 in~\cite{Shafarevich}).
\begin{fact}\label{fact:dim-tangent}
	For any irreducible variety $\V{V}$ and any $p \in \V{V}$ we have
	$\dim \T_p \V{V} \ge \dim \V{V}$.
\end{fact}


We will also need the following easy observation about the interplay between tangents and intersections.

\begin{proposition}\label{prop:tangent-cap}
	For any two varieties $\V{V}$ and $\V{W}$, and any $p \in \V{V} \cap \V{W}$, 
	$$\T_p (\V{V} \cap \V{W}) \sub \T_p\V{V} \,\cap\, \T_p\V{W}.$$
	In particular, if $\V{V} \sub \V{W}$ then $\T_p \V{V} \sub \T_p \V{W}$.
\end{proposition}
\begin{proof}
	We have $\I(\V{V}) \sub \I(\V{V}\cap \V{W})$ and $\I(\V{W}) \sub \I(\V{V}\cap \V{W})$. Therefore, by the definition of a tangent space, for any $p \in \V{V}\cap\V{W}$ we have $\T_p(\V{V}\cap\V{W}) \sub \T_p\V{V}$ and $\T_p(\V{V}\cap\V{W}) \sub \T_p\V{W}$,
	and thus also $\T_p(\V{V}\cap\V{W}) \sub  \T_p\V{V} \cap  \T_p\V{W}$, as claimed.
\end{proof}

\subsection{Slice rank of tangent spaces of determinantal varieties}

We henceforth denote by $\V{M}_r = \V{M}_r(\KK^{m\times n}) \sub \KK^{m \times n}$ the variety of matrices in $\KK^{m\times n}$ of rank at most $r$.
Note that $\V{M}_r$ is indeed a variety, as it is cut out by a finite set of polynomials: all $(r+1)\times(r+1)$ minors.
It is therefore referred to in the literature as a \emph{determinantal variety}.

The following crucial lemma shows that certain tangent spaces of the variety $\V{M}_r=\V{M}_{r}(\KK^{m \times n})$,
which are matrix spaces, 
have a small slice rank (recall Subsection~\ref{subsec:notation} for the terminology).

\begin{lemma}[Slice rank of tangents]\label{lemma:SR-Mr}
	The tangent space to $\V{M}_r=\V{M}_r(\KK^{m \times n})$, for any algebraically closed field $\KK$, at any matrix $A \in \V{M}_r$ with $\rank(A)=r$ satisfies
	$$\SR(\T_A \V{M}_r) \le 2r.$$
\end{lemma}
To prove Lemma~\ref{lemma:SR-Mr} will need the following result,
which explicitly describes the tangent space to $\V{M}_r$ at any matrix of rank exactly $r$. 
It can be deduced from Example~14.16 in~\cite{Harris}.
We prove it below for completeness.
\begin{proposition}[Tangents of determinantal varieties]\label{fact:tangent-matrices}
	The tangent space to $\V{M}_r=\V{M}_r(\KK^{m \times n})$, for any algebraically closed field $\KK$, at any matrix $A \in \V{M}_r$ with $\rank(A)=r$ is
	$$\T_A \V{M}_r = \{C A + A C' \mid C \in \KK^{m\times m}, C' \in \KK^{n \times n}\}.$$
\end{proposition}
\begin{proof}
	It will be convenient to work with the following equivalent definition of a tangent space of a variety $\V{V}$ at a point $p \in \V{V}$;
	$$\T_p\V{V} = \{ \B{v} \in \KK^n \mid \forall g \in \I(\V{V}) \colon g(p+t\B{v})-g(p) \equiv 0 \pmod{t^2}\}.$$
	To see this equivalence, observe that, using the Taylor expansion of the polynomial $g$ at the point $p$, we have 
	$g(p+t\B{v})-g(p) \equiv t\frac{\partial g}{\partial \B{v}}(p) \pmod{t^2}$.
	
	Now, we will use the fact that the $(r+1)\times(r+1)$ minors not only cut out the variety $\V{M}_r=\V{M}_r(\KK^{m\times n})$, but in fact generate the ideal $\I(\V{M}_r)$.
	Indeed, this follows from the fact that the ideal $I$ they generate is prime (\cite{BrunsVe}, Theorem 2.10) and so $\sqrt{I}=I$, together with Hilbert's Nullstellensatz which gives $\I(\V{M}_r)=\sqrt{I}=I$.
	Let $g_{I,J}$ denote the minor of the submatrix whose set of rows and columns are given by $I \sub [m]$ and $J \sub [n]$, respectively. 
	Thus, $\I(\V{M}_r) = \langle g_{I,J} \mid |I|=|J|=r+1 \rangle$.
	
	Since $\rank(A)=r$, there are invertible matrices $P \in \FF^{m\times m}$ and $Q \in \FF^{n \times n}$ such that $A=PI_rQ$, where 
	$$I_r = \begin{pmatrix}
		1 & 0 & \cdots & 0 & 0 & \cdots & \cdots & \cdots & 0 \\
		0 & 1 & \cdots & 0 & 0 & \cdots & \cdots & \cdots & 0 \\
		\vdots & \vdots & \ddots & \vdots & \vdots & \vdots & \vdots & \vdots & \vdots \\
		0 & 0 & \cdots & 1 & 0 & \cdots & \cdots & \cdots & 0 \\
		0 & 0 & \cdots & 0 & 0 & \cdots & \cdots & \cdots & 0\\
		\vdots  & \vdots & \vdots & \vdots & \vdots & \vdots & \vdots & \vdots & \vdots \\
		0 & 0 & \cdots & 0 & 0 & \cdots & \cdots & \cdots & 0
	\end{pmatrix}_{m \times n}$$
	has the $r \times r$ identity matrix as the upper-left submatrix, that is, the submatrix whose set of rows $I$ and set of columns $J$ are $I=[r]$ and $J=[r]$.
	
	Let $X \in \KK^{m \times n}$. Put $Y = P^{-1}XQ^{-1} \in \KK^{m \times n}$. 
	For every $g = g_{I,J}$ with $|I|=|J|=r+1$ 
	we have $g(A)=0$ and 
	$$g(A+tX) = g(P(I_r+tY)Q) = g(P)g(I_r+tY)g(Q).$$
	It follows that
	\begin{center}
		$g(A+tX)-g(A) \equiv 0 \pmod{t^2} \linebreak\quad\text{ if and only if }\quad 
		g(I_r+tY) \equiv 0 \pmod{t^2}$.
	\end{center}
	Write $Y=(y_{i,j})_{i,j}$. 
	Observe that if $I=[r]\cup\{i\}$ and $J=[r]\cup\{j\}$ for some $i>r$ and $j>r$ then $g_{I,J}(I_r+tY) \equiv ty_{i,j} \pmod{t^2}$, and otherwise $g_{I,J}(I_r+tY) \equiv 0 \pmod{t^2}$.
	Thus, $Y$ satisfies $g(I_r+tY) \equiv 0 \pmod{t^2}$ for every $g \in \I(\V{M}_r)$ if and only if $y_{i,j}=0$ for every $i>r$ and $j>r$, or equivalently, $Y=Y_1I_r + I_rY_2$ for some $Y_1 \in \KK^{m \times m}$ and $Y_2 \in \KK^{n \times n}$.
	We deduce 
	\begin{align*}
		\T_A\V{M}_r &= \{ X \in \KK^{m \times n} \mid \forall g \in \I(\V{M}_r) \colon\\
		& g(A+tX)-g(A) \equiv 0 \pmod{t^2}\}\\
		&= \{ PYQ \mid \exists Y_1 \in \KK^{m\times m}, Y_2 \in \KK^{n\times n} \colon Y = Y_1I_r + I_r Y_2\}\\
		&= \{ (PY_1P^{-1})A + A(Q^{-1} Y_2 Q) \mid Y_1 \in \KK^{m\times m}, Y_2 \in \KK^{n\times n}\}\\
		&= \{C A + A C' \mid C \in \KK^{m\times m}, C' \in \KK^{n \times n}\},
	\end{align*}
	completing the proof.
\end{proof}

We note that any matrix $A$ with $\rank(A)=r$ is a nonsingular point of $\V{M}_r$ (i.e., $\dim\T_A \V{M}_r = \dim\V{M}_r$), whereas any matrix $B$ with $\rank(B)<r$ is a singular point, and in fact, $\T_B\V{M}_r(\KK^{m \times n}) = \KK^{m \times n}$. 

\begin{proof}[Proof of Lemma~\ref{lemma:SR-Mr}]
	We identify $A=(a_{i,j}) \in \V{M}_r(\KK^{m \times n})$ with the bilinear form given by $A(\B{x},\B{y}) = \B{x}^T A \B{y} = \sum_{i,j} a_{i,j} x_i y_j$.
	Since $A \in \KK^{m \times n}$ and $\rank(A) \le r$, there are linear forms $f_1(\B{x}),\ldots,f_r(\B{x}) \in \KK[\B{x}]$ and linear forms $g_1(\B{y}),\ldots,g_r(\B{y}) \in \KK[\B{y}]$ such that
	$$A(\B{x},\B{y}) = \sum_{i=1}^r f_i(\B{x})g_i(\B{y}) .$$
	It follows that any matrix of the form $C A + A C'$, with $C \in \KK^{m\times m}$ and $C' \in \KK^{n \times n}$, has a corresponding bilinear form
	\begin{align}\label{eq:decomp}
		\begin{split}
			\B{x}^T (C A + A C') \B{y} 
			&= (C^T\B{x})^T A \B{y} + \B{x}^T A (C'\B{y})\\
			&= \sum_{i=1}^r f_i(C^T\B{x})g_i(\B{y}) + \sum_{i=1}^r f_i(\B{x})g_i(C'\B{y}) .
		\end{split}
	\end{align}
	Now, let $B_1,\ldots,B_d$ be any basis of $\T_A \V{M}_r$. 
	Then $\T_A \V{M}_r$ corresponds to the trilinear form $T=\sum_{k=1}^d z_k B_k(\B{x},\B{y})$ in the variables $\B{x},\B{y},\B{z}$. 
	By Proposition~\ref{fact:tangent-matrices}, 
	for each $k \in [d]$ we can write $B_k = C_kA+AC_k'$ for some $C_k \in \KK^{m\times m}$ and $C'_k \in \KK^{n \times n}$.
	Using the decomposition in~(\ref{eq:decomp}), we obtain the trilinear decomposition
	\begin{align*}
		T &= \sum_{k=1}^d z_k \cdot B_k(\B{x},\B{y})\\
		&= \sum_{k=1}^d z_k\Big(\sum_{i=1}^r f_i(C_k^T\B{x})g_i(\B{y}) + \sum_{i=1}^r f_i(\B{x})g_i(C_k'\B{y})\Big)\\
		&= \sum_{i=1}^r h_i(\B{x},\B{z}) g_i(\B{y}) + \sum_{i=1}^r f_i(\B{x})h'_i(\B{y},\B{z}) 
	\end{align*}
	where
	$$h_i(\B{x},\B{z}) := \sum_{k=1}^d z_kf_i(C_k^T\B{x}),\quad 
	h'_i(\B{y},\B{z}) := \sum_{k=1}^d z_k g_i(C_k'\B{y}) .$$
	Note that each $h_i \in \KK[\B{x},\B{z}]$ and $h'_i \in \KK[\B{y},\B{z}]$ are bilinear forms over $\KK$, 
	and recall that each $f_i \in \KK[\B{x}]$ and $g_i \in \KK[\B{y}]$ are linear forms over $\KK$.
	We deduce that each of the $2r$ summands in the decomposition of $T$ above is a trilinear form of slice rank at most $1$ over $\KK$.
	This completes the proof.
\end{proof}







\section{Slice rank vs. geometric rank}\label{sec:SR-GR}

In this section we prove the core of our main result, linearly bounding the slice rank of a tensor from above by its geometric rank.

\begin{theorem}\label{theo:main}
	For any $3$-tensor $T$ over any perfect field $\FF$,
	$$\SR(T) \le 3\GR(T).$$
\end{theorem}

We in fact get the slightly better constant $2$ instead of $3$ in Theorem~\ref{theo:main}, at the price of allowing the slice rank decomposition to use coefficients from an algebraic extension.

Let $\overline{\SR}(T)$ denote, for a tensor $T$, the slice rank over the algebraic closure of the field of coefficients of $T$. In other words, if $T$ is a tensor over $\FF$ then $\overline{\SR}(T)$ allows coefficients from the algebraic closure $\overline{\FF}$, rather than just from $\FF$, in the decomposition of $T$ into slice-rank one summands.
Clearly, $\overline{\SR}(T) \le \SR(T)$.
We note that for matrices, $\rank$ and $\overline{\rank}$ are equal. 
For tensors we have the following inequality, essentially due to Derksen~\cite{Derksen20} (we include a proof sketch at the end of this section).


\begin{proposition}[\cite{Derksen20}]\label{prop:stable-rk}
For any $3$-tensor $T$ over any perfect field,
$\frac23\SR(T) \le \overline{\SR}(T)$.
\end{proposition}




We will also need the following properties of slice rank, which are easily deduced from definition. 
For convenience of application, we state them for matrix spaces.

\begin{proposition}\label{prop:SR-properties}
The slice rank satisfies the following properties, where $\V{L}$ and $\V{L'}$ are linear subspaces of matrices:
\begin{enumerate}
	\item(Dimension bound) $\SR(\V{L}) \le \dim \V{L}$,
	\item(Monotonicity) $\SR(\V{L'}) \le \SR(\V{L})$ if $\V{L'} \preceq \V{L}$, 
	\item(Sub-additivity) $\SR(\V{L}+\V{L'}) \le \SR(\V{L})+\SR(\V{L'})$.
\end{enumerate}
\end{proposition}







\subsection{Linear sections of determinantal varieties}

For $\V{L} \preceq \KK^{m \times n}$ a matrix space we define the variety $\V{L}_r = \V{L} \cap \V{M}_{r}$ (here $\V{M}_{r}=\V{M}_{r}(\KK^{m \times n})$)
of all matrices in $\V{L}$ of rank at most $r$.
We next bound the slice rank of a matrix space using these linear sections of a determinantal variety. 
We denote by $\codim_L \V{X}$ the codimension of a variety $\V{X} \sub L$ inside a linear space $L$; that is, $\codim_L \V{X} = \dim L - \dim \V{X}$.

\begin{proposition}\label{theo:core}
Let $\V{L} \preceq \KK^{m \times n}$ be a matrix space over any algebraically closed field $\KK$. 
For any $r \in \NN$, 
$$\SR(\V{L}) \le 2r + \codim_\V{L} \V{L}_r.$$
\end{proposition}
\begin{proof}
We proceed by induction on $r$. Note that the base case $r=0$, which reads $\SR(\V{L}) \le 0+\codim_\V{L} \{\B{0}\} = \dim \V{L}$, follows from Proposition~\ref{prop:SR-properties}.
We thus move to the inductive step. 

Let $\V{V}$ be an irreducible component of $\V{L}_r$ with $\dim \V{V} = \dim \V{L}_r$,
and let $A \in \V{V} \sm \V{M}_{r-1}$.
We may indeed assume $\V{V} \sm \V{M}_{r-1} \neq \emptyset$, as otherwise $\V{V} \sub \V{L}_{r-1}$ and thus 
$\dim\V{L}_r = \dim\V{V} \le \dim \V{L}_{r-1}$ 
and we are done via the induction hypothesis by taking codimensions.
Let $\V{P} \preceq \V{L}$ be the linear subspace 
$\V{P} = \V{L} \,\cap\, \T_A\V{M}_r$. 
We will prove:
\begin{enumerate}
	\item\label{item:SR-bd} $\SR(\V{P}) \le 2r$, 
	\item\label{item:codim-bd} $\codim_\V{L} \V{P} \le \codim_\V{L} \V{L}_r$.
\end{enumerate}
To see why this would complete the inductive step, 
let $\V{P}^\perp$ be a complement subspace of $\V{P}$ in $\V{L}$, 
and note that
$$\SR(\V{L}) \le \SR(\V{P}) + \SR(\V{P}^\perp) \le \SR(\V{P}) + \codim_{\V{L}} \V{P} \le 2r + \codim_{\V{L}} \V{L}_r$$
where the first and second inequalities use Proposition~\ref{prop:SR-properties}, and the third inequality uses Items~(\ref{item:SR-bd}) and~(\ref{item:codim-bd}). 

For the proof of Item~(\ref{item:SR-bd}), we have
$$\SR(\V{P}) = \SR(\V{L} \cap \T_A \V{M}_r) \le \SR(\T_A \V{M}_r) \le 2r$$
where the first inequality uses Proposition~\ref{prop:SR-properties}, and the second inequality uses Lemma~\ref{lemma:SR-Mr} as $\rank(A)=r$.
For the proof of Item~(\ref{item:codim-bd}), we have
\begin{align*}
	\dim\V{L}_r = \dim \V{V} &\le \dim \T_A \V{V} \\ 
	&\le \dim \T_A \V{L}_r\\
	&\le \dim( \T_A \V{L} \,\cap\, \T_A \V{M}_r) \\
	&= \dim(\V{L} \,\cap\, \T_A \V{M}_r) 
	= \dim \V{P}
\end{align*}
where the first inequality uses Fact~\ref{fact:dim-tangent}, 
the second inequality uses Proposition~\ref{prop:tangent-cap} together with the fact that $\V{V} \sub \V{L}_r$, 
the third inequality uses Proposition~\ref{prop:tangent-cap} again,
and the last equality uses $\T_A \V{L}=\V{L}$ since $\V{L}$ is a linear subspace.
As the above varieties are subvarieties of $\V{L}$, we obtain
$\codim_{\V{L}} \V{P} \le \codim_{\V{L}} \V{L}_r$.
This proves Item~(\ref{item:codim-bd}) and therefore completes the proof of the inductive step.
\end{proof}

We note that an immediate corollary of Proposition~\ref{theo:core}
is a slice rank upper bound of $2r$ for any subspace of matrices of rank at most $r$.

\subsection{Putting everything together}

To prove Theorem~\ref{theo:main} we also need the following characterization of geometric rank.
Recall that $\GR(T)=\codim \ker T$ where $\ker T=\{(\B{x},\B{y}) \mid T(\B{x},\B{y},\cdot)=\B{0} \}$.
\begin{fact}[\cite{KoppartyMoZu20}]\label{fact:GR}
For any $3$-tensor $T$ over any field,
$$\GR(T) = \min_r \, r + \codim\{\B{x} \mid \rank T(\B{x},\cdot,\cdot)=r\} .$$
\end{fact}
Fact~\ref{fact:GR} is proved via the decomposition
$$\ker T = \bigcup_r \{(\B{x},\B{y}) \in \ker T \mid \rank T(\B{x},\cdot,\cdot)=r\} ,$$
using a result from algebraic geometry on the dimensions of fibers, and the fact that the codimension of a finite union of varieties is the minimum of their codimensions.
We refer to Theorem~3.1 in~\cite{KoppartyMoZu20} for the formal proof.

We are now ready to prove the main result of this section.
First, we show how to obtain Proposition~\ref{prop:stable-rk} from the results in~\cite{Derksen20}.


\begin{proof}[Proof sketch of Proposition~\ref{prop:stable-rk}]
This is obtained by combining Theorem~2.5, Corollary~3.7, and Proposition~4.9 in~\cite{Derksen20}.
These results show that the ``$G$-stable rank'' $\rank^G_\FF(T)$ over a perfect field $\FF$ satisfies the following properties, respectively: 
\begin{itemize}
	\item $\rank^G_\FF(T)=\rank^G_{\overline{\FF}}(T)$,
	\item $\rank^G_\FF(T) \le \SR(T)$,
	\item $\rank^G_\FF(T) \ge (2/3)\SR(T)$.
\end{itemize}
Putting these together gives
$\frac23 \SR(T) \le \rank^G_\FF(T) = \rank^G_{\overline{\FF}}(T) \le \overline{\SR}(T)$, as claimed.
\end{proof}

\begin{proof}[Proof of Theorem~\ref{theo:main}]
Suppose $T=(a_{i,j,k})_{i,j,k} \in \FF^{n_1 \times n_2 \times n_3}$ with $\FF$ an arbitrary field.
Let $\V{L} \preceq \overline{\FF}^{n_2\times n_3}$ be the matrix space spanned by the $n_1$ slices $A_1=(a_{1,j,k})_{j,k},\ldots,A_{n_1}=(a_{n_1,j,k})_{j,k}$. 
Note that we may assume, by acting with general linear group $\GL_{n_1}(\FF)$ on $T$, that the first $d:=\dim\V{L}$ slices $A_1,\ldots,A_{d}$ of $T$ are linearly independent and the rest are zero matrices; indeed, this action does not change $\GR(T)$ (see Lemma~4.2 in~\cite{KoppartyMoZu20}) nor does it change $\SR(T)$.

Note that for any $\B{x} \in \overline{\FF}^{n_1}$, the bilinear form $T(\B{x},\cdot,\cdot)$ corresponds to the matrix $\sum_i x_iA_i$;
indeed, 
\begin{align*}
	T(\B{x},\cdot,\cdot) &\colon (\B{y},\B{z}) \mapsto \sum_{i,j,k} a_{i,j,k}x_iy_jz_k = \sum_i x_i \sum_{j,k} a_{i,j,k}y_jz_k \\
	&= \sum_i x_i \B{y}^TA_i\B{z} = \B{y}^T\Big(\sum_i x_iA_i\Big)\B{z}.
\end{align*}
Using our assumption that $A_i = \B{0}$ for every $i>d$, 
let 
\begin{align*}
	\V{X}_r &= \{\B{x} \in \overline{\FF}^{n_1} \mid \rank T(\B{x},\cdot,\cdot) \le r\}\\ 
	&= \{\B{x} \in \overline{\FF}^{n_1} \mid \rank\big(x_1A_1+\cdots+x_dA_d\big) \le r\}.
\end{align*}
We claim that $\codim\V{X}_r=\codim_\V{L} \V{L}_r$. 
Recall that $\V{L}_r = \{ A \in \V{L} \mid \rank A \le r\}$. 
First, we show that the variety $\V{X}_r$ is isomorphic to the variety $\V{L}_r \times \overline{\FF}^{n_1-d}$. 
Indeed, the polynomial map (in fact linear) 
$$(x_1,\ldots,x_{n_1}) \mapsto (x_1A_1+\cdots+x_dA_d,x_{d+1},\ldots,x_{n_1})$$ 
maps $\V{X}_r$ to $\V{L}_r\times \overline{\FF}^{n_1-d}$, and is invertible via a polynomial map (in fact linear) by our assumption that $A_1,\ldots,A_{d}$ are linearly independent.
We deduce from this isomorphism the equality of dimensions $\dim\V{X}_r = \dim(\V{L}_r\times \overline{\FF}^{n_1-d})$, 
or equivalently, $\codim\V{X}_r = n_1-\dim\V{X}_r = d-\dim\V{L}_r = \codim_\V{L} \V{L}_r$, as claimed.


Let $r$ achieve the minimum in Fact~\ref{fact:GR}. This implies that
$\GR(T) = r + \codim \V{X}_r$. 
By Theorem~\ref{theo:core},
\begin{align*}
	\overline{\SR}(T) &= \SR(\V{L}) \le 2r + \codim_\V{L} \V{L}_r
	= 2r + \codim \V{X}_r\\
	&= 2\GR(T) - \codim \V{X}_r 
	\le 2\GR(T).
\end{align*}
Assuming further that $\FF$ is a perfect field and using Proposition~\ref{prop:stable-rk}, 
we finally obtain the bound
$\SR(T) \le \frac32\overline{\SR}(T) \le 3\GR(T)$, 
as desired.
\end{proof}




\section{Geometric rank vs.\ analytic rank}\label{sec:GR-AR}

Our main result in this section gives an essentially tight upper bound on the geometric rank in terms of the analytic rank. 

\begin{proposition}\label{prop:GR-AR}
For any $3$-tensor $T$ over any finite field $\FF$,
$$\AR(T) \ge (1-\log_{|\FF|}2)\GR(T) .$$
\end{proposition}



\subsection{Schwartz-Zippel meet B\'{e}zout}

We will need a certain generalized version of the classical Schwartz-Zippel lemma that applies to varieties.
We note that there are various generalized versions of the Schwartz-Zippel lemma appearing in the literature (e.g., Lemma~14 in~\cite{BukhTs12}, Claim~7.2 in \cite{DvirKoLo14}, Lemma A.3 in~\cite{EllenbergObTa10}).
However, in our version below the bound goes down exponentially with the codimension of the variety as soon as the field is larger than the degrees of the polynomials cutting out the variety, which is crucial for proving Proposition~\ref{prop:GR-AR}.

We use the notation $\V{V}(\FF) := \V{V} \cap \FF^n$ for any variety $\V{V} \sub \overline{\FF}^n$ defined over $\FF$.
Recall that a variety $\V{V} = \VV(f_1,\ldots,f_m)$ is said to be cut out by the polynomials $f_1,\ldots,f_m$.


\begin{lemma}[Schwartz-Zippel for varieties]\label{lemma:SZ}
Let $\FF$ be a finite field.
For any variety $\V{V} \sub \overline{\FF}^n$
cut out by polynomials of degrees at most $d$, 
$$\frac{|\V{V}(\FF)|}{|\FF|^n} \le \Big(\frac{d}{|\FF|}\Big)^{\codim\V{V}}.$$
\end{lemma}

We note that the classical Schwartz-Zippel lemma is recovered as the special case of Lemma~\ref{lemma:SZ} where $\V{V}$ is cut out by a single polynomial $p$. Indeed, in this case, Lemma~\ref{lemma:SZ} says that if $p$ is a non-zero polynomial, meaning $\codim\V{V}=1$, then $|\V{V}(\FF)|/|\FF|^n \le d/|\FF|$.

Let $\V{V}^0$ denote the union of the $0$-dimensional irreducible components of a variety $\V{V}$. Note that $\V{V}^0$ is a finite set.
For the proof of Lemma~\ref{lemma:SZ} we will use the overdetermined case of B\'{e}zout's inequality, which provides an upper bound on $|\V{V}^0|$ (see~\cite{Tao11}, Theorem 5).

\begin{fact}[Overdetermined B\'{e}zout's inequality]\label{fact:bezout-over}
Let $\V{V} = \VV(f_1,\ldots,f_m) \sub \KK^n$ be a variety, for an algebraically closed field $\KK$, 
cut out by $m \ge n$ polynomials. 
Write $\deg f_1 \ge \cdots \ge \deg f_m \ge 1$. 
Then 
$$|\V{V}^0| \le \prod_{i=1}^n \deg f_i.$$
\end{fact}

The \emph{degree} of an equidimensional\footnote{All irreducible components have the same dimension.} variety $\V{V} \sub \KK^n$, denoted $\deg\V{V}$, is the cardinality of the intersection of $\V{V}$ with a generic linear subspace in $\KK^n$ of dimension $\codim\V{V}$ (a well-defined, finite number).
The degree of an arbitrary variety $\V{V}$ is the sum of the degrees of its irreducible components.
The proof of Lemma~\ref{lemma:SZ} will ``bootstrap'' the following generalization of the Schwartz-Zippel lemma. 

\begin{fact}[\cite{BukhTs12},\cite{DvirKoLo14}]\label{fact:SZ}
Let $\FF$ be a finite field.
For any variety $\V{V}$ over $\overline{\FF}$,
$$|\V{V}(\FF)| \le \deg\V{V}\cdot|\FF|^{\dim\V{V}}.$$
\end{fact}


\begin{proof}[Proof of Lemma~\ref{lemma:SZ}]
We claim that the following inequality holds
assuming $\V{V}$ is equidimensional;
$$\deg \V{V} \le d^{\codim \V{V}}.$$
Suppose $\V{V}$ is cut out by $m$ polynomials of degree at most $d$. 
Note that $m \ge \codim\V{V}$.
Consider the variety obtained by intersecting $\V{V}$ with a generic linear subspace in $\overline{\FF}^n$ of dimension $\codim\V{V}$, 
and observe that it can be embedded as a variety $\V{W} \sub \overline{\FF}^{n_0}$ with $n_0=\codim\V{V}$.
Then $\V{W}$ satisfies the following properties:
\begin{itemize}
	\item $\dim\V{W} = 0$,
	\item $\V{W}$ is cut out by $m$ polynomials of degree at most $d$.
\end{itemize}	
In particular, and similarly to the above, $m \ge n_0$.
It follows that 
$$\deg \V{V} = |\V{W}| = |\V{W}^0| \le d^{n_0} = d^{\codim\V{V}},$$
where the first equality is by the definition of $\deg\V{V}$, the second equality uses $\V{W}=\V{W}^0$ as $\dim\V{W}=0$, 
and the inequality applies Fact~\ref{fact:bezout-over} since we are in the overdetermined case $m \ge n_0$.

To finish the proof, we need to remove the assumption in the above inequality that $\V{V}$ is equidimensional.
This is immediate using the standard technique of replacing the polynomials cutting out $\V{V}$ with generic linear combinations thereof, giving a variety $\V{V'} \supseteq \V{V}$ over $\overline{\FF}$ with $\dim\V{V'}=\dim\V{V}$ that is an equidimensional set-theoretic complete intersection (see, e.g., Kollar~\cite{Kollar88}, Theorem~1.5).\footnote{Alternatively, this can be done by defining an appropriate variant of degree for varieties that are not equidimensional.}
Indeed, 
starting with a single polynomial (which trivially cuts out an equidimensional variety), 
adding each such linear combination $p$, 
that is, intersecting with the hypersurface $\V{H}_p$ corresponding to $p$, reduces by one the dimension of every irreducible component $\V{X}$
(since $\V{X} \nsubseteq \V{H}_p$ as long as $\V{X} \nsubseteq \V{V}$).
Now, apply Fact~\ref{fact:SZ} to obtain
$$\frac{|\V{V}(\FF)|}{|\FF|^n} 
\le \frac{|\V{V'}(\FF)|}{|\FF|^n} 
\le \frac{d^{\codim\V{V}}|\FF|^{\dim\V{V}}}{|\FF|^n}
= \Big(\frac{d}{|\FF|}\Big)^{\codim\V{V}},$$
as desired.
\end{proof}


\subsection{Putting everything together}

We now deduce the desired bound relating $\GR$ and $\AR$.
We will use the following well known characterization of $\AR$; we include a proof for completeness.
\begin{fact}\label{fact:AR-ch}
For any $3$-tensor $T \in \FF^{n_1 \times n_2 \times n_3}$ over any finite field $\FF$,
$$\AR(T) = -\log_{|\FF|} \Pr_{\B{x},\B{y}}[f(\B{x},\B{y})=\B{0}]$$
where $f \colon \FF^{n_1} \times \FF^{n_2} \to \FF^{n_3}$ is the bilinear map corresponding to $T$.
\end{fact}
\begin{proof}
Write $f=(f_1,\ldots,f_{n_3})$, so that $T(\B{x},\B{y},\B{z})=\sum_{k=1}^{n_3} f_k(\B{x},\B{y})z_k$.
Denote $\bias(T) = \Exp_{\B{x},\B{y},\B{z}} \chi(T(\B{x},\B{y},\B{z}))$,
where $\chi$ is an arbitrary, nontrivial additive character of $\FF$,
so that $\AR(T) = -\log_{|\FF|}\bias(T)$.
Since $f$ is bilinear, we have $\bias(T) = 
\Pr_{\B{x},\B{y}} [f(\B{x},\B{y})=\B{0}]$; 
indeed, 
\begin{align*}
	\bias(T) &= \Exp_{\B{x},\B{y},\B{z}} \chi(T(\B{x},\B{y},\B{z}))
	= \Exp_{\B{x},\B{y},\B{z}} \chi\Big(\sum_k f_k(\B{x},\B{y})z_k \Big)\\
	&= \Exp_{\B{x},\B{y}} \Exp_{\B{z}}\prod_k\chi(f_k(\B{x},\B{y}) z_k)
	= \Exp_{\B{x},\B{y}} \prod_k \Exp_{z \in \FF}\chi(f_k(\B{x},\B{y}) z)\\
	&= \Exp_{\B{x},\B{y}} \prod_k [f_k(\B{x},\B{y})=0]
	= \Exp_{\B{x},\B{y}} [f(\B{x},\B{y})=\B{0}]\\
	&= \Pr_{\B{x},\B{y}} [f(\B{x},\B{y})=\B{0}] ,
\end{align*}
where $[\cdot]$ is the Iverson bracket. 
This completes the proof.
\end{proof}

\begin{proof}[Proof of Proposition~\ref{prop:GR-AR}]
Suppose $T \in \FF^{n_1 \times n_2 \times n_3}$. 
Put $\V{V} = \ker(T) \sub \overline{\FF}^{N}$ with $N=n_1+n_2$.
By Lemma~\ref{lemma:SZ},
$|\V{V}(\FF)|/|\FF|^{N} \le (2/|\FF|)^{\codim\V{V}}$.
Using Fact~\ref{fact:AR-ch}, it follows that
$$\AR(T) = -\log_{|\FF|} \frac{|\V{V}(\FF)|}{|\FF|^{N}}
\ge \codim\V{V} \cdot (1-\log_{|\FF|}2) .$$
As $\GR(T) = \codim\V{V}$, we are done.
\end{proof}



We are finally ready to combine our various bounds and obtain the main result.

\begin{proof}[Proof of Theorem~\ref{main:summary}]
The first inequality is given by Theorem~\ref{theo:main}.
The second inequality follows from Proposition~\ref{prop:GR-AR} for any finite $\FF \neq \FF_2$, since
\begin{align*}
	\GR(T) &\le (1-\log_{|\FF|}2)^{-1}\AR(T)\\ 
	&\le (1-\log_{3}2)^{-1}\AR(T) \le 2.71\AR(T).
\end{align*}
\end{proof}
We note that, as evident from the proof of Theorem~\ref{main:summary}, we in fact obtain the bounds $\SR(T) \le 3\GR(T) \le 3(1+o_{|\FF|}(1))\AR(T)$.
